% Transcription — fonds Alexandre Grothendieck, shelfmark 151, batch 2 (pages 21-40).
% Established from the facsimile archives/batches/151/batch-02.pdf (a mirror of
% https://grothendieck.umontpellier.fr/151.pdf), per the skill
% transcribe-grothendieck.
%
% Page 25 is blank — absent. Every other page of the batch carries the
% mathematics and is transcribed.
%
% The batch holds two pieces of writing. Pages 21-24 are numbered 4 to 7 by the
% author in a circle at the head of the page: they continue, without a break,
% the series numbered 1 to 3 that closed batch 1 (pages 16, 18, 19), and the
% coordinate calculation of page 20 — d(x,y) = (x,y,y), rho^2 d(x,y) = (y,x,y) —
% is resumed on page 21 in the same terms. Pages 26-40 are the text
% « Structures stratifiées » itself, whose table of contents opened batch 1 on
% page 1; it runs continuously, and the author numbers his sheets 1 to 8 at the
% head of alternate archivist pages (26, 28, 30, 32, 34, 36, 38, 40), each sheet
% being written on both sides. The batch stops mid-sentence at the foot of
% page 40 — the text goes on into the next batch.
%
% Two diagrams, (8) and the first (23), draw the two arrows that close the
% square as dotted lines: those arrows are what the van Kampen assertion is to
% supply, and drawing them solid would assert what the page is asking. The
% renderer's tikz-cd subset was extended with `dashed` in the same commit.
%
% On uncertainty: \uncertain{} marks a reading that carries mathematical weight
% and that we do not hold to be sure; \ill marks what we do not read at all.
% Function words whose sense is forced by the sentence — a « où », a « donc » —
% are given plainly. The pages numbered 4 to 7 are working sheets, written over,
% boxed and cancelled; there the notes say where we have compressed.
%
% The dating is the inventory's own and spans the group's range; the notes
% themselves carry no date.
%
% Pass: Opus 5 (claude-opus-5), 2026-08-15 — first pass, unchecked against the pages by a human.
\documentclass[11pt,a4paper]{article}
% Le préambule commun à toutes les transcriptions du fonds Grothendieck.
%
% Il tient en une idée : **l'incertitude se compose, elle ne se résout pas.**
% Une transcription qui lisse un mot douteux a détruit la seule chose pour
% laquelle on la faisait — un lecteur doit pouvoir savoir, à chaque endroit,
% ce qui était sur le feuillet et ce qui a été ajouté. Les six macros qui
% suivent sont donc l'appareil critique, et non de la décoration.
%
% Les mêmes commandes sont reconnues par `scripts/render.mjs`, qui produit la
% vue de lecture du site. Une macro ajoutée ici doit l'être aussi là-bas,
% faute de quoi le rendu échouera bruyamment — ce qui est voulu : un rendu
% silencieusement faux est pire qu'un rendu absent.

\ProvidesPackage{grothendieck}[2026/08/08 Transcriptions du fonds Grothendieck]

\usepackage{amsmath,amssymb,amsthm}
\usepackage{mathrsfs}
\usepackage{stmaryrd}
% Les diagrammes commutatifs sont partout dans ce fonds : c'est la langue
% même de la géométrie algébrique de Grothendieck, et une transcription qui
% les rendrait en prose ne transcrirait rien.
\usepackage{tikz-cd}
% Français par défaut — c'est la langue des feuillets — mais l'anglais est
% chargé aussi : la traduction et le résumé sont des documents anglais, et la
% typographie française (espaces avant les deux-points) y serait une faute.
% Ces documents basculent par \selectlanguage{english} en tête de corps.
\usepackage[english,french]{babel}
\usepackage{fontspec}
\usepackage{geometry}
\geometry{a4paper,margin=25mm,marginparwidth=28mm}
\usepackage{marginnote}
\usepackage{xcolor}
% ulem, not soul. `soul`'s \st and \ul are built on a character-by-character
% scan that XeLaTeX's font machinery breaks: the document fails with an
% undefined control sequence rather than degrading. ulem does the same two jobs
% and is Unicode-engine safe. `normalem` stops it hijacking \emph, which the
% transcriptions use for Grothendieck's own emphasis.
\usepackage[normalem]{ulem}
\usepackage{hyperref}
\hypersetup{colorlinks=true,linkcolor=black,urlcolor=black,pdfborder={0 0 0}}

% --- Métadonnées du lot -----------------------------------------------------
% Reprises telles quelles par le script de rendu, qui les affiche en tête de
% la vue de lecture. Elles fixent ce que le fichier prétend couvrir : sans
% elles, une transcription flotte sans point d'attache dans un fonds de seize
% mille feuillets.

\newcommand{\folder}[1]{\gdef\grothFolder{#1}}
\newcommand{\batch}[1]{\gdef\grothBatch{#1}}
\newcommand{\leaves}[2]{\gdef\grothFirst{#1}\gdef\grothLast{#2}}
\newcommand{\foldertitle}[1]{\gdef\grothFoldertitle{#1}}
\gdef\grothFolder{?}\gdef\grothBatch{?}\gdef\grothFirst{?}\gdef\grothLast{?}\gdef\grothFoldertitle{}

% --- Le résumé --------------------------------------------------------------
%
% Il ouvre la lecture modernisée, sous le titre « Résumé », et il est composé
% en petit. Ce n'est pas un ornement : c'est le seul passage du document qui
% ne soit pas des mathématiques, et un lecteur doit pouvoir le reconnaître
% avant de l'avoir lu — puis le sauter s'il connaît déjà le sujet, ou s'y
% arrêter s'il ne le connaît pas. Le filet qui le suit marque la reprise du
% corps.

\newenvironment{resume}
  {\small\setlength{\parskip}{0.45em}}
  {\par\medskip\hrule\medskip}

% --- Mots-clés ---------------------------------------------------------------
%
% En anglais, en clôture du résumé de la lecture modernisée : le vocabulaire
% moderne sous lequel chercher ce que les feuillets construisent. Le manifeste
% du site (`npm run manifest`) les extrait pour étiqueter la cote entière —
% cette ligne est la source unique des tags, il n'y a pas de fichier à côté.
\newcommand{\keywords}[1]{\par\smallskip\noindent
  {\footnotesize\textsc{Keywords} --- \textit{#1}}\par}

% --- L'appareil critique ----------------------------------------------------

% Le numéro de feuillet, en marge. C'est l'ancre qui permet de régler un
% désaccord sur une formule en regardant *un* feuillet plutôt qu'une chemise
% de six cents. Le site s'en sert aussi pour tourner les pages du fac-similé
% quand on fait défiler la transcription.
\newcommand{\leaf}[1]{%
  \marginnote{\footnotesize\color{gray!70}\textsf{#1}}%
  \label{leaf:#1}\ignorespaces}

% Illisible : on ne devine pas. Un mot inventé qui se lit comme les autres est
% le pire résultat possible.
\newcommand{\ill}{\textcolor{red!60!black}{[\,\ldots\,]}}

% Lecture douteuse : le mot est donné, et signalé comme incertain.
\newcommand{\uncertain}[1]{\uline{#1}}

% Ajout de l'éditeur — une lettre manquante, un mot sous-entendu.
\newcommand{\add}[1]{\textcolor{blue!50!black}{[#1]}}

% Biffé par Grothendieck lui-même. Ce qu'il a rayé fait partie du document :
% une rature dit souvent où la pensée a changé de direction.
\newcommand{\struck}[1]{\sout{#1}}

% Note du transcripteur, sur l'état matériel du feuillet.
\newcommand{\note}[1]{\par\smallskip{\small\color{gray!60!black}\textit{[#1]}}\par\smallskip}

% Note marginale de Grothendieck, distincte du corps du texte.
\newcommand{\marginal}[1]{\par\smallskip\noindent\hspace{1em}%
  {\small\color{gray!60!black}#1}\par\smallskip}

% --- Titre ------------------------------------------------------------------

\AtBeginDocument{%
  \begin{center}
    {\large\bfseries Fonds Alexandre Grothendieck --- cote n\textsuperscript{o}~\grothFolder}\\[2pt]
    {\itshape\grothFoldertitle}\\[4pt]
    {\small feuillets \grothFirst--\grothLast\ (lot \grothBatch)}\\[2pt]
    {\footnotesize\color{gray!60!black}Transcription établie d'après le fac-similé numérisé
     par l'Université de Montpellier.\\
     Les lectures incertaines sont soulignées, les illisibles notées [\,\ldots\,],
     les ajouts entre crochets.}
  \end{center}
  \vspace{1em}}

\endinput


\folder{151}
\batch{2}
\pages{21}{40}
\dating{[à partir de 1981-à partir de 1990]}
\watermark{Édition de démonstration}
\foldertitle{Topos stratifié (1981) : notes manuscrites (s.d.).}

\begin{document}

\section*{Le système $X_1, X_2, X_3$ — suite}

\note{Les quatre pages qui suivent sont numérotées 4 à 7 par l'auteur, dans un
cercle en haut de page : elles poursuivent sans coupure la série numérotée 1 à
3 du lot précédent. Ce sont des feuilles de travail, où les relations sont
écrites, encadrées, biffées et reprises. Nous transcrivons les systèmes de
relations et les phrases sans reproduire chaque reprise ; là où un fragment est
recouvert au point de ne plus se lire, il est noté.}

\page{21}
\[
  \mathfrak{S}_2 = \{1,\sigma\}, \qquad
  \mathfrak{S}_3 = \{1,\rho,\rho^2,\sigma_1,\sigma_2,\sigma_3\},
\]
\[
  \sigma_i^2 = 1, \qquad
  \bigl[\ \rho^3 = 1, \quad
  \rho\,\sigma_i\,\rho^{-1} = \sigma_{i+1} \quad (\sigma_4 = \sigma_1)\ \bigr].
\]

\begin{tikzcd}[column sep=large]
X_1 & X_2 \arrow[r, "d", bend left=18] & X_3 \arrow[l, "p", bend left=18]
\end{tikzcd}

avec $X_1 = \{e_1\}$, $e_2 \in X_2$, $e_3 \in X_3$, et $\sigma e_2 = e_2$.

Les opérations sont données en coordonnées :

\[
  d(x,y) = (x,y,y), \qquad
  \rho\,d(x,y) = (y,y,x), \qquad
  \rho^2 d(x,y) = (y,x,y),
\]
\[
  p(x,y,z) = (x,y),
\]
\[
  \rho(x,y,z) = (y,z,x), \qquad \rho^2(x,y,z) = (z,x,y),
\]
\[
  \sigma_1(x,y,z) = (x,z,y), \qquad
  \sigma_2(x,y,z) = (z,y,x), \qquad
  \sigma_3(x,y,z) = (y,x,z).
\]

Les relations encadrées :

\[
  \boxed{p\,d = \mathrm{id}_{X_2}}, \qquad
  \boxed{p\,\rho\,d = e}, \qquad
  \boxed{p\,\rho^2 d = \sigma},
\]
\[
  \boxed{\sigma_1 d = d}, \qquad \boxed{p\,\sigma_3 = \sigma\,p},
\]

où $e$ est l'application constante de valeur $e_2$.

\note{La seconde relation encadrée ne s'accorde pas aux formules en
coordonnées portées plus haut sur la même page : $p\rho d(x,y) = (y,y)$, qui
n'est pas constant. Nous transcrivons l'une et l'autre telles quelles ; la
page ne tranche pas.}

\[
  \mathfrak{S}_3 = \{1,\sigma_1\}\cdot\{1,\rho,\rho^2\}
                 = \{1,\rho,\rho^2\}\cdot\{1,\sigma_1\}.
\]

$d(X_2) \subset X_3$, et l'on cherche quand deux éléments $\rho^i d(\xi)$ et
$\rho^{i'} d(\xi')$ coïncident : cela revient à $i \equiv i' \pmod 3$ et
$\xi = \xi'$, ou bien $\xi = \xi' = e_2$. Posant $d(\xi') = \rho^j d(\xi)$ :

\[
  \begin{cases}
    \text{si } j = 0 : & \xi' = \xi \\
    \text{si } j = 1 : & \xi' = p\,\rho\,d(\xi) = e_2 \\
    \text{si } j = 2 : & \xi' = p\,\rho^2 d(\xi) = \sigma\xi
  \end{cases}
\]

et dans le cas $j = 1$, $d(\xi') = d(e_2) = e_3$, donc
$\xi = p\,d(\xi) = p\,e_3 = e_2$, donc $\xi = \xi' = e_2$. En regard :
$d\sigma(\xi) = \rho^2 d(\xi)$, $d(\xi) = \rho\,d\sigma(\xi)$, d'où
$\xi = p\,\rho\,d\sigma(\xi) = e_2$ et $\xi' = e_2$.

\note{Cette discussion de cas occupe la moitié basse de la page, écrite deux
fois, la seconde par-dessus la première, avec plusieurs lignes biffées ; nous
donnons la chaîne complète et laissons les reprises. Un exposant de $\rho$ y
est raturé au point de ne plus se lire.}

\page{22}
Ainsi, le sous-$\mathfrak{S}_3$-ens. $X'_3$ de $X_3$ engendré par $d(X_2)$ a
une unique orbite à 1 élément, savoir $\{e_3\} = \{d\,e_2\}$, toutes les autres
orbites sont d'ordre 3 (dans le \ill), et elles sont en corr. biunivoque avec
les éléments de $X_2 \setminus \{e_2\} = X^{*}_2$. Comme
$\mathfrak{S}_3$-ens. $X'_3$ s'écrit conséquemment

\[
  X'_3 \simeq \{e_3\} \amalg
    \bigl(\mathfrak{S}_3 \wedge_{\{1,\sigma_1\}} X^{*}_2\bigr),
\]

où $\{1,\sigma_1\}$ opère trivialement sur $X^{*}_2$. On pose

\[
  X^{*}_3 \overset{\text{déf}}{=} X_3 \setminus
    \underbrace{X'_3}_{\mathfrak{S}_3\cdot d(X_2)}
\]

comme $\mathfrak{S}_3$-ens., et l'application induite par $p$

\begin{tikzcd}[column sep=large]
X^{*}_3 \arrow[r, "p^{*} = p|X^{*}_3"] & X_2
\end{tikzcd}

satisfont à : $\boxed{p^{*}\sigma_3 = \sigma\,p^{*}}$.

\marginal{\ill}

\note{Le bas de la page gauche porte, écrit en diagonale sur cinq lignes, un
paragraphe où reviennent $X'_3$ et un renvoi $(*)$ ; il est trop recouvert
pour être lu et nous ne le devinons pas.}

\[
  X_3 \xrightarrow{\ (p,\,p\rho)\ } X_2 \times X_2 \quad \text{épi} :
\]

pour $\xi, \eta \in X_2$, $\exists\,\alpha \in X_3$ avec $p\alpha = \xi$,
$p\rho\alpha = \eta$, où l'on note $p(x,y,z) = (x,y)$ et
$p'(x,y,z) = (y,z) = p\rho(x,y,z)$.

\page{23}
« Quelle est la partie de $X_2 \times X_2$ qui provient de $X'_3$ ? »

\begin{enumerate}
  \item[a)] $(p,p\rho)(e_3) = (e_2,e_2)$
  \item[b)] $\displaystyle (p,p\rho)\,\rho^i d(\xi) =
    \begin{cases}
      (\xi, e_2) & \text{si } i = 0 \\
      (e_2, \sigma\xi) & \text{si } i = 1 \\
      (\sigma\xi, \xi) & \text{si } i = 2
    \end{cases}$
    \quad $i \in \mathbb{Z}/3\mathbb{Z}$, $\xi \in X^{*}_2 = X_2\setminus\{e_2\}$.
\end{enumerate}

Donc on trouve les éléments de

\[
  \underbrace{(X_2 \times \{e_2\})}_{\Sigma_1} \cup
  \underbrace{(\{e_2\} \times X_2)}_{\Sigma_2} \cup \Delta'
\]

où $\Delta'$ est l'image de $X_2 \xrightarrow{\ \delta'\ } X_2 \times X_2$,
$\delta' = (\sigma, \mathrm{id}_{X_2})$.

NB Cette partie $\Sigma$ de $X_2 \times X_2$ est la réunion des $\Sigma_1$,
$\Sigma_2$, $\Delta'$, dont l'intersection deux à deux (et prise des trois) est
égale à $\{(e_2,e_2)\}$.

Quand est-ce qu'on a $\Sigma = X_2 \times X_2$ ? On a
$X_2 \times X_2 \setminus (\Sigma_1 \cup \Sigma_2) = X^{*}_2 \times X^{*}_2$,
donc il est \ill\ que $X^{*}_2 \to X^{*}_2 \times X^{*}_2$ soit surj., i.e. que
pour $\xi, \eta \in X^{*}_2$, on ait tjs $\eta = \sigma\xi$. Cela signifie
qu'on est dans l'un des 2 cas suivants :

\begin{enumerate}
  \item[a)] $X^{*}_2 = \varnothing$ i.e. $X_2 = \{e_2\}$
  \item[b)] $X^{*}_2$ de cardinal 1 (donc $\sigma$ y opère trivialement)
    \marginal{i.e. $X_2$ de cardinal 2 (et $\sigma$ y opère trivialement)}
\end{enumerate}

[Si $X^{*}_2$ est de cardinal $\geq 2$, alors s'il y a une seule orbite sous
$\sigma$, celle-ci est de cardinal 2, et pour $\xi \in X^{*}_2$,
$(\xi,\xi) \notin \mathrm{Im}\,\delta'$ i.e. $\xi \neq \sigma\xi$ ; ou s'il y a
au moins deux orbites, on prend $\xi, \eta$ dans deux orbites différentes et on
a bien $\eta \neq \sigma\xi$.]

\page{24}
Supposons donc $X_2$ de cardinal 2, et $X_3 = X'_3$
\quad ($X_2 = \{e_2, \xi\}$), donc

\[
  X_3 = \{e_3\} \amalg \{d\xi,\ \rho\,d\xi,\ \rho^2 d\xi\}.
\]

Considérons les deux carrés, avec
$\delta_1 = d = \bigl((x,y) \mapsto (x,y,y)\bigr)$ et
$\delta_3 = \bigl((x,y) \mapsto (x,x,y)\bigr)$ :

\begin{tikzcd}[column sep=small, row sep=small]
X_3 \arrow[r, "\delta_1"] & X_2 \arrow[d, "\delta_0"] \\
X_2 \arrow[u, "\delta_3"] \arrow[r, "\delta_0"'] & X_1 = \{e_1\}
\end{tikzcd}
\qquad
\begin{tikzcd}[column sep=small, row sep=small]
X_2 & X_3 \arrow[l, "p = \mathrm{pr}_2"'] \\
X_1 = \{e_1\} \arrow[u, "\kappa_1"] & X_2 \arrow[l, "p_1"] \arrow[u, "\kappa'_2"']
\end{tikzcd}

sont-ils cartésiens ? On a $\delta_3(x,y) = (x,x,y)$, et
$d\sigma(x,y) = d(y,x) = (y,x,x) = \rho^2\delta_3(x,y)$, i.e.

\[
  \delta_3 = \rho^{\ill}\,d\sigma .
\]

\note{L'exposant de $\rho$ est raturé et ne se lit plus. Les carrés sont
étiquetés $(*)$ et $(**)$ dans la marge gauche.}

Dans un cas où l'on suppose seulement $\mathrm{card}\,X_1 = 1$ i.e.
$X_1 = \{e_1\}$ : quand a-t-on, pour $\xi, \xi' \in X_2$,
$d(\xi) = \rho^{\ill}\,d\sigma(\xi')$ ? Si $\xi' = e_2$, alors
$\xi = \xi'' = e_2$ et c'est bien le cas. Reste la question

\[
  p^{-1}(e_2) \overset{?}{=} \delta_3(X_2)
    \quad \bigl[\ = \rho\,d\sigma(X_2) = \rho\,d(X_2)\ \bigr],
\]

l'inclusion $\supset$ étant acquise a priori, et $X_3 = X'_3
\overset{\text{déf}}{=} \mathfrak{S}_3 \cdot d(X_2)$.

$2^{\circ}$\ \ $\xi \in X'_3 \setminus \{e_3\}$ : il s'agit de prouver que
$\xi \in \rho\,d(X_2)$. On a $p(\xi) = e_2$, et $\xi = \rho^i d(\eta)$ avec
$\eta \in X^{*}_2$ et $i \in \{0,1,2\}$, et

\[
  p\xi = \begin{cases}
    \eta \ (\neq e_2) & \text{si } i = 0 \\
    \sigma\eta \ (\neq e_2) & \text{si } i = 2
  \end{cases}
\]

donc $p\xi = e_2 \implies i = 1$, \emph{qed}.

\note{La moitié basse de cette page est un enchaînement d'essais raturés qui
reprennent trois fois la même vérification ; nous n'en gardons que la chaîne
qui aboutit, celle que l'auteur clôt par « qed ».}

\section*{Structures stratifiées}

\note{Ce qui suit est le texte annoncé par le sommaire ouvrant le lot
précédent. Il est d'une seule venue, et l'auteur numérote ses feuillets 1 à 8
en haut de page : le numéro paraît une page sur deux, le verso poursuivant le
recto sans le porter.}

\page{26}
\note{Feuillet 1 de l'auteur.}

\textbf{Structures stratifiées.}

1. La situation la plus élémentaire, en un sens qui apparaîtra, sera la
suivante. On se donne des espaces (disons ! \ill\ en attendant qu'on se décide
sur les topos \ldots)

\[
  (1) \qquad Y \hookrightarrow X
\]

$Y$ fermé dans $X$. On suppose que la « structure normale » de $X$ le long de
$Y$ soit loc. triviale, i.e. que pour tout $y \in Y$, existe un voisinage
ouvert $X'$ de $y$ et un espace $Z$ pointé par $z$, et un diagramme commutatif
(2) (où $Y' = Y \cap X'$)

\begin{tikzcd}[column sep=large]
Y' \arrow[r, hook] \arrow[d, "\mathrm{can}\,\wr"'] & X' \arrow[d, "\wr"] \\
Y'\times\{z\} \arrow[r, hook, "\mathrm{id}_{Y'} \times i_{z,Z}"'] & Y'\times Z
\end{tikzcd}

où $X' \simeq Y'\times Z$ est un homéomorphisme. On considère alors le
voisinage tubulaire (« infinitésimal ») de $Y$ dans $X$, soit

\[
  (3) \qquad Y \hookrightarrow \mathcal{V}_{Y,X} \hookrightarrow X ,
\]

dont il faudra donner une définition \page{27} techniquement souple.

Si $X$ est \struck{un espace} un espace paracompact, on pourra prendre (à titre
provisoire !) pour $\mathcal{V}_{Y,X}$ le « germe d'espace » de $X$ autour de
$Y$ — qui est \ill\ un topos, sav. la $\varprojlim$ topologique des voisinages
ouverts de $Y$ dans $X$.

Intuitivement, $\mathcal{V}_{Y,X} \hookrightarrow X$ a les propriétés d'un
ouvert, et $Y \hookrightarrow \mathcal{V}_{Y,X}$ celles d'une immersion fermée,
et (3) constitue une factorisation canonique de (1). L'hypothèse de locale
trivialité normale implique (moyennant des conditions de « modération »
ensemblistes) une équivalence d'homotopie

\[
  (4) \qquad Y \xrightarrow{\ \text{équiv. d'homotopie}\ } \mathcal{V}_{Y,X}.
\]

\marginal{L'intuition \ill\ $\mathcal{V}_{Y,X}$ et $Y$ \ill\ équivalents \ill\
(indépendamment de \ill\ de locale trivialité).}

Supposons \ill\ $Y$ et $X - Y = X^{*}$ \struck{des} variétés, et
$Z^{*} = Z \setminus \{z\}$ une variété (moyennant des hyp. de modération,
l'hyp. sur $Z$ devrait être une conséquence de la première hypothèse), alors
que $z$ est un pt singulier de $Z$, i.e. $X$ est « singulier le long de $Y$ ».
Il est naturel alors de vouloir reconstituer $X$ en termes des variétés $Y$ et
$X^{*}$, \page{28} jouant le rôle de « pièces de construction »
élémentaires. Ceci amène à introduire

\[
  (5) \qquad \mathcal{V}^{*}_{Y,X} \overset{\text{déf}}{=}
    \mathcal{V}_{Y,X} \cap X^{*} = \mathcal{V}_{Y,X} \setminus Y
    \qquad (X^{*} = X \setminus Y),
\]

et le diagramme (6)

\begin{tikzcd}[column sep=large]
\mathcal{V}^{*}_{Y,X} \arrow[r, "\text{bouchage de trous le long de } Y"] \arrow[d] & \mathcal{V}_{Y,X} \arrow[d] \\
X^{*} \arrow[r] & X
\end{tikzcd}

Je vais regarder (indépendamment de toute hypothèse — de locale trivialité
normale, — de lissité de $X^{*}$ et $Y$) $X$ comme obtenu par
\emph{recollement} de $\mathcal{V}_{Y,X}$ et de $X^{*}$ suivant leur
intersection $\mathcal{V}^{*}_{Y,X}$, ou (en termes plus techniques) comme la
somme amalgamée, comme espace ou topos \uncertain{limite} du diagramme (7)

\begin{tikzcd}[column sep=small, row sep=small]
\mathcal{V}^{*}_{Y,X} \arrow[r, hook] \arrow[d] & \mathcal{V}_{Y,X} \\
X^{*} &
\end{tikzcd}

\page{29}
Si on dispose d'une théorie suffisamment souple des types d'homotopie, on
devrait pouvoir en déduire le type d'homotopie de $X$ comme « somme
amalgamée » des types d'homotopie des trois \struck{types} « \ill\ »
$\mathcal{V}^{*}_{Y,X}$, $\mathcal{V}_{Y,X}$ et $X^{*}$ (dont les deux premiers
sont des êtres « infinitésimaux »). Par exemple, on doit avoir pour les espaces
déduits en « tuant les $\Pi_i$ pour $i \geq 2$ » — donc, essentiellement, pour
les groupoïdes fondamentaux, un « théorème de van Kampen », qui s'exprime par
le fait que le diagramme des groupoïdes fondamentaux (8)

\begin{tikzcd}[column sep=large, nodes={font=\scriptsize}]
\Pi_1\mathcal{V}^{*}_{Y,X} \arrow[r, "p"] \arrow[d, "i"'] & \Pi_1\mathcal{V}_{Y,X}\ (\simeq \Pi_1 Y) \arrow[d, dashed] \\
\Pi_1 X^{*} \arrow[r, dashed] & (\Pi_1 X)
\end{tikzcd}

\note{Les deux flèches qui ferment ce carré sont pointillées sur la page, et
les deux termes de droite mis entre parenthèses : c'est ce que l'énoncé de van
Kampen doit fournir, non ce qu'on tient déjà.}

\page{30}
\note{Feuillet 3 de l'auteur.}

est cocartésien — ou encore, si $Y$, $X$, $X^{*}$ sont connexes, et qu'on
choisisse un « pt base » de $\Pi_1\mathcal{V}^{*}_{Y,X}$ (i.e. par définition,
un rev. universel de $\mathcal{V}^{*}_{Y,X}$) (d'où pts bases de
$\Pi_1\mathcal{V}_{Y,X}$, $\Pi_1 X^{*}$, $\Pi_1 X$), qu'on a un isom. can. des
groupes fondamentaux

\[
  (9) \qquad \Pi_1(X) \simeq \Pi_1(X^{*})
    \wedge_{\Pi_1\mathcal{V}^{*}_{Y,X}}
    \underbrace{\Pi_1(\mathcal{V}_{Y,X})}_{\simeq\ \Pi_1(Y)} ,
\]

où $\Pi_1(\mathcal{V}_{Y,X})$ est isom. extérieurement à $\Pi_1(Y)$.

Pour expliciter $\Pi_1(X)$ en termes des données « élémentaires » $\Pi_1(Y)$ et
$\Pi_1(X^{*})$, il \uncertain{vaudrait} \struck{encore} à expliciter la
structure de $\Pi_1(\mathcal{V}^{*}_{Y,X})$, qui s'envoie dans l'un et dans
l'autre, donnant lieu justement au diagramme (10)

\begin{tikzcd}[column sep=small, row sep=small]
\Pi_1(\mathcal{V}^{*}_{Y,X}) \arrow[r] \arrow[d] & \Pi_1(Y) \\
\Pi_1(X^{*}) &
\end{tikzcd}

\marginal{cette flèche définie mod. des choix \ill\ (pt base de $Y$ \ill).}

qui exprime (8). C'est là que l'hypothèse \page{31} de
\emph{locale trivialité normale} a un rôle à jouer (celle de lissité \ill\
devenant techniquement inutile, son rôle était plutôt de nature heuristique
\ldots). On doit se débrouiller, dans ce cas, pour démontrer que les propriétés
homotopiques de l'inclusion
$\mathcal{V}^{*}_{Y,X} \hookrightarrow \mathcal{V}_{Y,X}$ sont celles d'une
\emph{fibration localement triviale de fibre $\mathcal{V}^{*}_{z,Z}$} :

\[
  (11) \qquad \mathcal{V}^{*}_{Y,X} \longrightarrow \mathcal{V}_{Y,X}
\]

est une « fibration homotopique » de fibre $\mathcal{V}^{*}_{z,X}$, de base
$Y \xrightarrow{\ \simeq\ \text{homot.}\ } \mathcal{V}_{Y,X}$ — et c'est cela
qui devrait servir, dans le contexte topologique (p. ex. celui des schémas sur
le typ. étale) de \emph{définition} de la « locale trivialité normale » de
l'inclusion $Y \hookrightarrow X$.

\note{La fibre est écrite $\mathcal{V}^{*}_{z,Z}$ dans la phrase soulignée et
$\mathcal{V}^{*}_{z,X}$ dans (11), deux lignes plus bas. Nous donnons chaque
occurrence telle qu'elle est portée.}

(Bien sûr, dans le contexte schématique, il faudra de plus travailler avec des
\struck{types} \ill\ types d'homotopie profinis, et même « localiser » ces types
d'homotopie \page{32} en l'un des nbs premiers qui sont distincts des
caractéristiques résiduelles qui interviennent, \ill\ que dans ce contexte
modifié on peut prouver des théorèmes qu'il faut \struck{de trivialité}, cf
Artin-Mazur \ldots).

On aura en particulier une suite exacte d'homotopie \struck{\ill}
(choisissant un « pt base » dans un sens \ill\ de la fibre \ill\
$\mathcal{V}^{*}_{z,Z}$) :

\[
  (12) \qquad \cdots \to \Pi_i(\mathcal{V}^{*}_{z,Z})
    \to \Pi_i(\mathcal{V}^{*}_{Y,X}) \to \Pi_i(Y)
    \to \Pi_{i-1}(\mathcal{V}^{*}_{z,Z}) \to \cdots ,
\]

qui nous intéressera surtout en basses dimensions (13)

\begin{tikzcd}[column sep=small, row sep=small, nodes={font=\scriptsize}]
\Pi_2(Y) \arrow[r] & \Pi_1(\mathcal{V}^{*}_{z,Z}) \arrow[r] & \Pi_1(\mathcal{V}^{*}_{Y,X}) \arrow[r] & \Pi_1(Y) \arrow[r] & \Pi_0(\mathcal{V}^{*}_{z,Z}) \arrow[d] \\
& & & & \Pi_0(\mathcal{V}^{*}_{Y,X}) \arrow[d] \\
& & & & \Pi_0(Y)
\end{tikzcd}

\note{Une accolade réunit sous (13) les termes
$\Pi_1(\mathcal{V}^{*}_{Y,X}) \to \Pi_1(Y)$, avec la mention « qui intervient
dans (10) ».}

Si on suppose p. ex. que $z$ est \uncertain{déconnectant} vers $Z$ localement,
i.e. que la « fibre » $\mathcal{V}^{*}_{z,Z}$ est connexe (ou encore, que $Y$
est \uncertain{déconnectant} dans $X$ (loc.)), alors
$\Pi_1(\mathcal{V}^{*}_{Y,X})$ est une \emph{extension} de $\Pi_1(Y)$ par un
quotient de $\Pi_1(\mathcal{V}^{*}_{z,Z})$ \page{33} (savoir le quotient
par l'image de $\Pi_2(Y)$), et $=$ par $\Pi_1(\mathcal{V}^{*}_{z,Z})$
lui-même si $\Pi_2(Y)$ est nul.

Si on suppose pour simplifier que
$\Pi_2(Y) \to \Pi_1(\mathcal{V}^{*}_{z,Z})$ est nul, on trouve donc un
diagramme (14) qui précise (10)

\begin{tikzcd}[column sep=small, row sep=small, nodes={font=\scriptsize}]
1 \arrow[r] & \Pi_1(\mathcal{V}^{*}_{z,Z}) \arrow[r] & \Pi_1(\mathcal{V}^{*}_{Y,X}) \arrow[r] \arrow[d] & \Pi_1(Y) \arrow[r] & 1 \\
& & \Pi_1(X^{*}) & &
\end{tikzcd}

où la ligne horizontale est exacte.

\note{Deux noms sont portés sous cette suite : « ($I$ ``inertie'') » sous
$\Pi_1(\mathcal{V}^{*}_{z,Z})$, et « ($D$ ``décomposition'') » à droite de la
flèche verticale. Nous rapportons la place qu'ils occupent sans décider à quel
terme le second se rapporte.}

Un cas intéressant est celui où les « pièces de construction » qui
interviennent, savoir $Y$, $X^{*}$ et $\mathcal{V}^{*}_{z,Z}$, sont des espaces
« $K(\pi,1)$ », i.e.

\[
  (15) \qquad \Pi_i(Y),\ \Pi_i(X^{*}),\ \Pi_i(\mathcal{V}^{*}_{z,Z})
    \text{ nuls pour } i \geq 2 ,
\]

alors on en déduit qu'il en est de même pour $\mathcal{V}^{*}_{Y,X}$ (et bien
sûr pour $\mathcal{V}_{Y,X} \simeq_{\text{homot.}} Y$).

Si on a la « théorie souple des \emph{liens} des types d'homotopie » à laquelle
je faisais \page{34} \note{Feuillet 5 de l'auteur.} allusion, on devrait
pouvoir exprimer alors \emph{le type d'homotopie} de $X$ (et non seulement les
$\Pi_i$) en termes du diagramme des groupoïdes 8, — ou, ce qui revient au
même, du diagramme de groupes (10).

En tout cas, il est clair (indépendamment de toute hypothèse de nullité de
certains $\Pi_i$, ou de locale trivialité normale) comment reconstruire, en
termes du diagramme 8, la catégorie $\mathfrak{F}$ dans la cat. des faisceaux
sur $X$, \struck{dont} formée des $F$ tels que l'on ait

\[
  (16) \qquad F|X^{*} \text{ et } F|Y \text{ localement triviaux} ,
\]

Cette catégorie $\mathfrak{F}$ est équivalente en effet à celle des systèmes

\[
  (17) \qquad (E_{X^{*}},\ E_{Y,X}\,;\ \varphi)
\]

En regard de (17), et biffé : \struck{(en un rev. étale de $X^{*}$)}, le mot
« étale » étant réécrit au-dessus.

où $E_{X^{*}}$ est un système local sur $\Pi_1 X^{*}$, $E_{Y,X}$ un syst. local
sur $\Pi_1\mathcal{V}_{Y,X}$ ($\simeq \Pi_1 Y$) (donc un rev. ét. de $Y$) et
enfin un hom. de syst. locaux sur $\Pi_1\mathcal{V}^{*}_{Y,X}$

\page{35}
\[
  (18) \qquad \varphi \colon p^{*}(E_{Y,X}) \longrightarrow i^{*}(E_{X^{*}}) .
\]

En termes des diagrammes de groupes

(1\ill) Quand les trois groupoïdes composants de (8) sont connexes, i.e. $Y$,
$X^{*}$ et $\mathcal{V}^{*}_{Y,X}$ connexes, on peut aussi redéfinir en
considérant $E_{X^{*}}$ comme un ens. sur lequel $\Pi_1(X^{*})$ opère,
$E_{Y,X}$ un ens. sur lequel $\Pi_1(Y)$ opère, et $p^{*}(E_{Y,X})$ et
$i^{*}(E_{X^{*}})$ dans (18) comme étant ces mêmes ensembles, comme
$\Pi_1(\mathcal{V}^{*}_{Y,X})$-ensembles par restriction des opérateurs.

\marginal{NB dans \uncertain{(14)} \ill}

\note{Une note de sept lignes, écrite en travers dans la marge gauche et
numérotée (18 bis), y reprend $E_{Y,X}$, $E_{X^{*}}$, $\Pi_1(Y)$ et $I$ ; elle
est trop recouverte pour être lue au-delà de ces termes.}

Un cas particulièrement typique pour la géométrie algébrique, est celui où $X$
est une surface top. et où $Y$ est réduit à un pt $a$, donc
$\mathcal{V}_{Y,a} = D_a$, \emph{disque} (infinitésimal) (ouvert) de $a$, et
$\mathcal{V}^{*}_{Y,a} = D^{*}_a$, disque épointé (inf.) autour de $a$, de
sorte qu'on trouve le diagramme \page{36} (19)

\begin{tikzcd}[column sep=small, row sep=small]
D^{*}_a \arrow[r, hook, "\text{« bouchage de trou en } a\text{ »}"] \arrow[d] & D_a \\
X^{*} = X \setminus \{a\} &
\end{tikzcd}

Ici le diagramme (10) se réduit à (20)

\begin{tikzcd}[column sep=small, row sep=small]
\mathbb{Z}_{a,X} \arrow[r] \arrow[d] & \{1\} \\
\Pi_1(X^{*}) &
\end{tikzcd}

où $\mathbb{Z}_{a,X}$ ($\simeq \mathbb{Z}$) est le « module d'orientation » de
$X$ en $a$, dont les deux gén. correspondent aux deux orientations de $X$ en
$a$. L'homom. vertical de (20) est l'homom. qui définit la structure « locale
en $a$ » du groupe fondamental $X^{*}$. Le cas le plus simple est celui où

\[
  (21) \qquad \begin{cases} X = \mathbb{C} \text{ plan complexe} \\ a = \{0\} \end{cases}
\]

d'où

\begin{tikzcd}[column sep=small, row sep=small]
D^{*}_0 \arrow[r] \arrow[d] & D_0 \\
\mathbb{C}^{*} &
\end{tikzcd}

dont le produit avec lui-même un nb arbitraire de fois, donne les
« modèles-types » des stratifications associées aux diviseurs à croisements
normaux (dans les \page{37} domaines analytiques complexes).

Ici $X = \mathbb{C}$ est topologiquement trivial, ce qui peut s'exprimer ainsi
par le fait que son type d'homotopie est complètement décrit par celui des
types $\mathfrak{F}$ des systèmes (17). En effet, on aura

\[
  (22) \qquad \mathfrak{F} \simeq \widehat{C}
\]

où $C$ est le cas qui est \marginal{où dans le diagramme (19)
$D \Rightarrow \Pi_1(\mathcal{V}^{*}_{Y,X}) \simeq \Pi_1(X^{*})$ iso}
défini comme la catégorie dont l'un des objets est l'unique \ill\ de (18), et
dont les flèches sont, en plus des flèches de ces trois groupoïdes, les
flèches d'un objet $a$ de $A = \Pi_1\mathcal{V}^{*}_{Y,X}$ vers un objet $b$
de $B = \Pi_1(\mathcal{V}_{Y,X})$ ou de $C = \Pi_1(X^{*})$, définies comme
étant les flèches

\[
  p^{*}(a) \longrightarrow b \ \text{ dans } B
  \qquad (\text{resp. } i^{*}(a) \longrightarrow c \ \text{ dans } C).
\]

\note{Les deux premiers tiers de cette page sont recouverts de longues ratures
horizontales et de deux traits obliques ; le texte que nous donnons est ce qui
subsiste sous elles, et la phrase qui définit $C$ y est reprise deux fois. Nous
ne reproduisons pas la première version, entièrement biffée.}

Dans le cas \struck{actuel} présent, les groupoïdes $A$, $B$, $C$ et les
morphismes entre eux \page{38} \note{Feuillet 7 de l'auteur.} où
$I = D$ est le groupe des auto. de $a$, et où les seules autres flèches, à part
les $\gamma \in D$ et l'identité de $b$, est une flèche $a \to b$ (de sorte que
$b$ est objet final de $C$). Comme $C$ admet un objet final, $\widehat{C}$ est
(comme aussi $\mathfrak{F}$) de type d'homotopie trivial, ce qui signifie bien
que $\mathfrak{F}$ — le type d'homotopie de \ill\ $X = \mathbb{C}$.

Prenons maintenant le cas où $X$ est le compactifié $\mathbb{P}^1_{\mathbb{C}}$
de $\mathbb{C}$ par un pt à l'infini, disons

\[
  \begin{cases}
    X = \mathbb{P}^1_{\mathbb{C}} \simeq S^2 \\
    Y = \{\infty\} \\
    X^{*} = \mathbb{C}
  \end{cases}
\]

le diagramme (7) devient (23)

\begin{tikzcd}[column sep=small, row sep=small]
D^{*}_{\infty} \arrow[r, hook] \arrow[d] & D_{\infty} \arrow[d, dashed] \\
\mathbb{C} \arrow[r, dashed] & \mathbb{P}^1_{\mathbb{C}} \simeq S^2
\end{tikzcd}

dont la somme amalgamée est $S^2$, dont le type d'homotopie n'a rien de
trivial ! Le \struck{type} $\mathfrak{F}$ se décrit maintenant à l'aide du
diagramme correspondant sur les $\Pi_1$, savoir \page{39} (23)

\begin{tikzcd}[column sep=small, row sep=small]
\mathbb{Z} \arrow[r] \arrow[d] & 1 \\
1 &
\end{tikzcd}

\note{L'auteur donne à ce diagramme le numéro (23), déjà porté par celui de la
page précédente.}

on a maintenant une équivalence de types

\[
  \mathfrak{F} \simeq \widehat{C} ,
\]

où $C$ est la catégorie ${\cdot} \to {\cdot}$, qui a un objet final — donc
$\mathfrak{F}$ est encore \emph{homotopiquement trivial}, alors que celle
\ill\ $X$ \emph{ne l'est pas} ! Pourtant les pièces de construction
$\{\infty\}$, $X^{*} \simeq \mathbb{C}$, et
$\mathcal{V}^{*}_{z,Z} \simeq D^{*}_{\infty}$ sont bel et bien des $K(\pi,1)$ —
(les deux premières sont contractiles !). Cela montre que le type
$\mathfrak{F}$, qui en principe semblait « tenir compte » des $\Pi_1$ des
« pièces de construction » $Y$, $X^{*}$, $\mathcal{V}^{*}_{Y,X}$ voisines, ne
donne pas le type d'homotopie de $X$, même si les « morceaux » sont des
$K(\pi,1)$ — en \uncertain{d'autres termes} le morphisme canonique

\[
  (24) \qquad X \longrightarrow \mathfrak{F}
\]

n'est pas une équivalence d'homotopie.

\page{40}
\note{Feuillet 8 de l'auteur.}

J'avais craint un moment, par cet exemple, que le type \struck{\ill}
d'homotopie de $X$ ne pouvait pas (sous les hypothèses $K(\pi,1)$ faites) se
récupérer à partir de ceux de $X^{*}$, $Y$, $\mathcal{V}^{*}_{Y,X}$ et du
diagramme (7), i.e. du diagramme (8) ou (10). Mais je me suis convaincu ensuite
que cette crainte n'était pas fondée, ou si : le type $\mathfrak{F}$ n'est pas
suffisant, lui, pour exprimer le type d'homotopie \ill.

Essayons en effet, dans le cas qui nous occupe — le cas de la division
cellulaire la plus simple de $S^2$, (qui est en même temps le cas le plus
\struck{simple} élémentaire d'un diviseur sur une courbe alg.
\uncertain{complexe} !) — de récupérer le $H^2(X, A)$ ($A$ un groupe abélien,
disons) à partir du diagramme (23). J'interprète $H^2(X,A)$ comme classifiant
les \emph{gerbes} sur $X$ liées par $A$ — et je constate que pour
\struck{de telles} gerbes, sur des types \ill\ une théorie des « \ill\ »

\note{Le texte s'arrête ici, au bas de la page, en milieu de phrase : il se
poursuit sur le feuillet suivant, hors de ce lot.}

\end{document}
